\documentclass[12pt]{article}
\usepackage[margin=1in]{geometry}
\usepackage{amsmath,amssymb}
\usepackage{graphicx}
\usepackage{booktabs}
\usepackage{hyperref}
\usepackage{natbib}
\usepackage{setspace}
\usepackage{lineno}
\doublespacing
\linenumbers

\title{Multi-Dimensional Neural and Behavioral Signatures of Acute Ketamine: Linking Inter-Individual Variability to Cortical Gene Expression Patterns of SST and PVALB Interneurons}

\author{Author One$^{1,*}$, Author Two$^{1}$, Author Three$^{2}$ \\
\small $^1$Department of Psychiatry, Yale University School of Medicine, New Haven, CT, USA \\
\small $^2$Department of Neuroscience, Yale University School of Medicine, New Haven, CT, USA \\
\small $^*$Corresponding author: author.one@yale.edu}

\date{}

\begin{document}
\maketitle

\begin{abstract}
Ketamine, an NMDA receptor antagonist with rapid-acting antidepressant properties, produces substantial inter-individual variability in neural and behavioral responses that remains poorly understood. Here we investigated whether ketamine's acute effects are uni-dimensional or multi-dimensional using a data-driven principal component analysis (PCA) framework applied to global brain connectivity (GBC) changes in 40 healthy participants undergoing a double-blind, placebo-controlled crossover study. We identified five significant principal components of ketamine-induced change in GBC (delta-GBC), collectively explaining 42.1\% of total variance, demonstrating that ketamine's neural effects are irreducibly multi-dimensional. Ketamine produced higher effective dimensionality (12.8 $\pm$ 0.7) compared to LSD (8.7 $\pm$ 0.3) and psilocybin (8.6 $\pm$ 0.3; $F(2,60)=564$, $p<0.001$). The principal neural gradient (PC1 delta-GBC) spatially correlated with cortical gene expression patterns of somatostatin (SST) and parvalbumin (PVALB) interneuron markers from the Allen Human Brain Atlas, implicating specific GABAergic interneuron subtypes in ketamine's acute effects. Behavioral symptoms were also multi-dimensional, with three significant behavioral PCs mapping onto distinct neural connectivity gradients at the single-subject level. These findings establish a multi-dimensional framework for understanding ketamine's pharmacological effects and support the feasibility of personalized pharmacological biomarkers.
\end{abstract}

\section{Introduction}

Ketamine, a non-competitive NMDA receptor antagonist, has emerged as one of the most promising treatments for treatment-resistant depression, producing rapid antidepressant effects within hours of administration \citep{berman2000antidepressant, zarate2006randomized}. Despite this therapeutic promise, approximately 35\% of patients fail to respond to a single ketamine infusion, highlighting the critical importance of understanding individual variability in the drug's effects \citep{abdallah2017ketamine}. The mechanisms underlying this variability remain largely unknown, representing a significant barrier to the development of precision medicine approaches for ketamine therapy.

At the systems level, ketamine produces widespread changes in functional connectivity as measured by resting-state functional magnetic resonance imaging (fMRI). Global brain connectivity (GBC), which quantifies the average functional connection strength of each brain region with all other regions, has been shown to be sensitive to ketamine's acute effects \citep{cole2012global, anticevic2015nmda}. However, most prior investigations of ketamine's neural effects have relied on seed-based connectivity approaches that require a priori selection of regions of interest, potentially missing distributed system-level effects. Furthermore, inconsistent findings in the literature---such as contradictory reports of thalamo-cortical connectivity changes---may stem from unexamined inter-individual variability rather than genuine population-level disagreement \citep{driesen2013relationship}.

A fundamental gap exists between ketamine's known molecular mechanism of action (NMDA receptor antagonism) and its diverse systems-level neural and behavioral effects. NMDA receptors are expressed throughout the brain, but different neuronal subtypes express distinct NMDA receptor subunit compositions with varying ketamine affinities \citep{tremblay2016gabaergic}. In particular, somatostatin-expressing (SST) and parvalbumin-expressing (PVALB) GABAergic interneurons preferentially express GluN2B- and GluN2A-containing NMDA receptors, respectively, suggesting that ketamine may differentially affect these cell populations. However, no framework has connected this molecular heterogeneity to the observed variability in systems-level connectivity changes.

At the behavioral level, ketamine produces a range of acute symptoms quantified by instruments such as the Positive and Negative Syndrome Scale (PANSS) \citep{kay1987positive} and cognitive tasks. The relationship between these behavioral effects and underlying neural connectivity changes has been examined primarily at the group level, obscuring potentially informative individual-level associations \citep{krystal1994subanesthetic}.

In this study, we addressed these gaps by applying a data-driven PCA framework to characterize the multi-dimensional structure of ketamine-induced GBC changes. We conducted a double-blind, placebo-controlled crossover study in 40 healthy participants, computing delta-GBC maps (ketamine minus placebo) across 718 functionally-defined brain parcels and performing PCA to identify the principal axes of neural variation. We compared the effective dimensionality of ketamine-induced changes with those of two classical psychedelics (LSD and psilocybin) using previously published datasets \citep{preller2019changes, carhart2016neural}. We then linked neural gradients to cortical gene expression patterns from the Allen Human Brain Atlas (AHBA) \citep{hawrylycz2012anatomically} and to multi-dimensional behavioral symptom variation at the single-subject level.

\section{Methods}

\subsection{Participants and Study Design}

Forty healthy adults participated in a double-blind, placebo-controlled crossover study. All participants were screened for psychiatric and neurological disorders and provided written informed consent. The study was approved by the local institutional review board and conducted in accordance with the Declaration of Helsinki.

\subsection{Drug Administration}

On two separate sessions, participants received either intravenous ketamine (initial bolus of 0.23 mg/kg followed by continuous infusion at 0.58 mg/kg/hour) or saline placebo, with session order counterbalanced across participants. Sessions were separated by at least one week to allow for complete drug washout.

\subsection{Neuroimaging Acquisition and Preprocessing}

Resting-state fMRI data were acquired during both ketamine and placebo sessions. Participants were instructed to rest with eyes open while maintaining fixation on a central crosshair. Standard preprocessing steps included motion correction, slice-timing correction, spatial normalization, and nuisance regression.

\subsection{Global Brain Connectivity Analysis}

The cortical surface was parcellated into 718 functionally-defined parcels. For each parcel, GBC was computed as the mean Fisher's $z$-transformed Pearson correlation of that parcel's time series with all other parcels. A delta-GBC map was created for each participant by subtracting the placebo GBC map from the ketamine GBC map. This yielded a $718 \times 40$ matrix of delta-GBC values for subsequent PCA.

\subsection{Principal Component Analysis of Neural Features}

PCA was performed on the delta-GBC matrix to identify the principal axes of inter-individual variation in ketamine-induced connectivity changes. The significance of each PC was assessed using a permutation test with 5,000 permutations ($p < 0.05$ threshold). The correlation between each significant PC map and the mean delta-GBC map was tested using Pearson correlation with Bonferroni correction ($p < 0.001$).

\subsection{Effective Dimensionality}

Effective dimensionality was computed as:
\begin{equation}
D_{\text{eff}} = \frac{\left(\sum_i \lambda_i\right)^2}{\sum_i \lambda_i^2}
\end{equation}
where $\lambda_i$ are the eigenvalues from PCA. This measure captures how many components are needed to account for the variance structure. Effective dimensionality was compared across ketamine (this study), LSD, and psilocybin (published datasets) using one-way ANOVA with post-hoc pairwise comparisons.

\subsection{Gene Expression Analysis}

Cortical gene expression patterns were obtained from the Allen Human Brain Atlas (AHBA) using the GEMINI-DOT framework. Expression maps from six postmortem brains were mapped onto the 718-parcel cortical surface. Spatial correlations were computed between the PC1 delta-GBC map and each gene's expression pattern across the approximately 17,000 available genes. Genes were ranked by their spatial correlation, and gene set enrichment analysis was performed to identify cell-type-specific signatures.

\subsection{Behavioral Assessment and PCA}

Behavioral effects of ketamine were assessed using the PANSS positive, negative, and general symptom subscales and a spatial working memory task. Behavioral change scores (ketamine minus placebo) were normalized and subjected to a separate PCA to identify principal axes of behavioral variation. The significance of behavioral PCs was assessed via permutation testing.

\subsection{Neuro-Behavioral Mapping}

For each participant, behavioral PC scores were regressed onto the delta-GBC map to generate neuro-behavioral maps. Regression coefficients were $Z$-scored across participants to identify brain regions where GBC changes reliably predicted behavioral symptom dimensions.

\section{Results}

\subsection{Multi-Dimensional Neural Effects of Ketamine}

PCA on the 718 $\times$ 40 delta-GBC matrix revealed five significant principal components (permutation test, $p < 0.05$), collectively capturing 42.1\% of total variance (Table~\ref{tab:pca_results}). PC1 and PC2 each explained more than 10\% of variance, while PC3--PC5 each accounted for approximately 5\%. Three of the five PCs (PC1, PC2, and PC4) were significantly correlated with the mean delta-GBC map (Bonferroni-corrected $p < 0.001$), while PC3 and PC5 captured patterns of variation orthogonal to the group average effect.

\begin{table}[ht]
\centering
\caption{Principal components of ketamine-induced delta-GBC.}
\label{tab:pca_results}
\begin{tabular}{lccl}
\toprule
Component & Variance Explained (\%) & Significance & Correlation with Mean \\
\midrule
PC1 & $>$10 & $p < 0.05$ & Significant ($p < 0.001$) \\
PC2 & $>$10 & $p < 0.05$ & Significant ($p < 0.001$) \\
PC3 & $\sim$5 & $p < 0.05$ & Not significant \\
PC4 & $\sim$5 & $p < 0.05$ & Significant ($p < 0.001$) \\
PC5 & $\sim$5 & $p < 0.05$ & Not significant \\
\midrule
Total (PCs 1--5) & 42.1 & -- & -- \\
\bottomrule
\end{tabular}
\end{table}

Critically, all five PCs were bi-directional: participants loading positively on a given PC showed the opposite pattern of GBC change from those loading negatively. For PC1, high positive loaders exhibited increased GBC in association networks (particularly the default mode network) and decreased GBC in sensory cortices (secondary visual cortex), while high negative loaders showed the reverse pattern. This bi-directionality provides a parsimonious explanation for contradictory findings in the prior ketamine resting-state literature, as the direction of connectivity change depends on where each individual falls along the principal gradient.

\subsection{Ketamine Shows Higher Effective Dimensionality Than Classical Psychedelics}

To contextualize the dimensionality of ketamine's neural effects, we compared effective dimensionality across ketamine, LSD, and psilocybin. As shown in Table~\ref{tab:dimensionality}, ketamine produced significantly higher effective dimensionality ($12.8 \pm 0.7$) than both LSD ($8.7 \pm 0.3$) and psilocybin ($8.6 \pm 0.3$; one-way ANOVA: $F(2,60) = 564$, $p < 0.001$). Post-hoc tests confirmed significant differences between ketamine and both psychedelics ($p < 0.001$), whereas LSD and psilocybin did not differ from each other ($p = 0.6$).

\begin{table}[ht]
\centering
\caption{Effective dimensionality comparison across pharmacological conditions.}
\label{tab:dimensionality}
\begin{tabular}{lccc}
\toprule
Substance & Dimensionality (mean $\pm$ SD) & vs.\ Ketamine & vs.\ LSD \\
\midrule
Ketamine & 12.8 $\pm$ 0.7 & -- & $p < 0.001$ \\
LSD & 8.7 $\pm$ 0.3 & $p < 0.001$ & -- \\
Psilocybin & 8.6 $\pm$ 0.3 & $p < 0.001$ & $p = 0.6$ \\
\midrule
\multicolumn{4}{l}{One-way ANOVA: $F(2,60) = 564$, $p < 0.001$} \\
\bottomrule
\end{tabular}
\end{table}

This result indicates that ketamine induces a substantially more complex pattern of inter-individual variation in brain connectivity than classical serotonergic psychedelics, potentially reflecting ketamine's broader receptor pharmacology or more heterogeneous downstream molecular targets.

\subsection{PC1 Delta-GBC Maps Onto SST and PVALB Interneuron Gene Expression}

To bridge the gap between systems-level connectivity changes and molecular mechanisms, we performed spatial correlation analysis between the PC1 delta-GBC map and cortical gene expression patterns from the AHBA. Across approximately 17,000 genes, the expression patterns of SST (somatostatin) and PVALB (parvalbumin) showed significant positive spatial correlations with the PC1 map (Table~\ref{tab:gene_expression}). Both of these genes are canonical markers for distinct subtypes of GABAergic interneurons that are known to preferentially express specific NMDA receptor subunits.

\begin{table}[ht]
\centering
\caption{Key gene expression correlations with PC1 delta-GBC map.}
\label{tab:gene_expression}
\begin{tabular}{llll}
\toprule
Gene / Class & Correlation Direction & Cell Type & NMDA Subunit \\
\midrule
SST & Significant positive & SST interneurons & GluN2B-containing \\
PVALB & Significant positive & PV interneurons & GluN2A-containing \\
Glutamate receptors & Varied & Multiple & Various \\
GABA receptors & Varied & Multiple & Various \\
\bottomrule
\end{tabular}
\end{table}

Gene set enrichment analysis confirmed that cell-type-specific gene signatures for SST and PVALB interneurons were significantly enriched among the genes whose cortical expression most closely tracked the principal gradient of ketamine-induced connectivity change. This finding provides a plausible link between ketamine's molecular action at NMDA receptors and the systems-level connectivity gradient captured by PC1.

\subsection{Multi-Dimensional Behavioral Effects}

All behavioral measures showed significant ketamine effects. PANSS positive, negative, and general symptom subscales were elevated under ketamine relative to placebo, while spatial working memory performance was impaired (all comparisons $p < 0.01$, Cohen's $d$ values indicated moderate-to-large effects).

Behavioral PCA identified three significant behavioral PCs (permutation testing, $p < 0.05$), indicating that the behavioral effects of ketamine are also multi-dimensional rather than reflecting a single symptom axis. Behavioral PC1 captured the largest share of variance with mixed loadings across PANSS subscales and cognition, while PC2 and PC3 represented dissociable dimensions of symptom variation. Individual variation in behavioral PC scores was substantial, paralleling the neural variability captured by the delta-GBC PCA.

\subsection{Neuro-Behavioral Mapping at the Single-Subject Level}

We next investigated whether the multi-dimensional behavioral variation could be mapped onto distinct neural connectivity gradients. Regression of behavioral PC1 scores onto individual delta-GBC maps revealed a robust spatial pattern of brain-behavior relationships, with association cortex regions (including default mode network components) showing positive relationships and sensory regions showing negative relationships with behavioral PC1 scores (Table~\ref{tab:neurobehavioral}).

\begin{table}[ht]
\centering
\caption{Properties of neuro-behavioral maps.}
\label{tab:neurobehavioral}
\begin{tabular}{lll}
\toprule
Property & PC1 Neuro-Behavioral Map & PC2 Neuro-Behavioral Map \\
\midrule
Association cortex & Positive relationship & Different pattern \\
Sensory cortex & Negative relationship & Different pattern \\
Thalamic involvement & Observed & Observed \\
Overlap with neural PCs & Partial & Partial \\
\bottomrule
\end{tabular}
\end{table}

The neuro-behavioral PC2 map revealed a distinct spatial pattern from PC1, confirming that different dimensions of behavioral variation are associated with separable neural substrates. Critically, these mappings were resolvable at the individual subject level: individual behavioral PC scores predicted individual delta-GBC spatial patterns, demonstrating that the neuro-behavioral framework captures meaningful individual differences rather than merely group-level associations.

Analysis of individual subjects confirmed wide variation in both behavioral and neuro-behavioral scores. Participants with high positive PC1 behavioral loadings showed increased GBC in association networks coupled with decreased GBC in sensory cortices, along with higher PANSS positive symptoms. Participants with high negative PC1 loadings showed the reverse neural pattern and lower symptom severity. Intermediate loaders showed mixed profiles, consistent with the continuous nature of the principal gradient.

\section{Discussion}

This study provides three major advances in understanding ketamine's acute effects. First, we demonstrate that ketamine-induced changes in brain connectivity are fundamentally multi-dimensional, with five significant PCs capturing 42.1\% of total variance. This finding challenges the implicit assumption in most pharmacological neuroimaging studies that a drug effect can be adequately characterized by a single group-average map. Second, we show that the principal neural gradient maps onto cortical gene expression patterns of SST and PVALB interneuron markers, providing a mechanistic bridge between ketamine's NMDA receptor antagonism and its systems-level connectivity effects. Third, we establish that behavioral symptom variation is also multi-dimensional and maps onto distinct neural connectivity gradients at the single-subject level, supporting the feasibility of personalized pharmacological biomarkers.

The finding that all five PCs are bi-directional has important implications for interpreting the existing ketamine literature. Prior studies reporting contradictory effects---such as both increased and decreased thalamo-cortical connectivity under ketamine---may have been observing the same multi-dimensional phenomenon from different vantage points, with sample composition determining which direction of the gradient dominated the group average. Our framework reconciles these discrepancies by showing that the direction of connectivity change depends on where each individual falls along the principal axes of variation.

The higher effective dimensionality of ketamine compared to LSD and psilocybin is noteworthy given their different pharmacological profiles. While classical psychedelics act primarily through serotonin 5-HT2A receptors, ketamine's primary target is the NMDA receptor, with additional actions at opioid, monoaminergic, and cholinergic systems. The broader pharmacological profile may contribute to the higher dimensionality by engaging multiple partially independent neural circuits. Alternatively, the greater dimensionality could reflect the fact that NMDA receptors are expressed by diverse neuronal populations with different subunit compositions, leading to a more heterogeneous pattern of effects across individuals.

The spatial correspondence between the PC1 delta-GBC map and SST/PVALB interneuron gene expression provides a cell-type-specific interpretation of the principal neural gradient. SST and PVALB interneurons are the two major classes of cortical GABAergic interneurons, and they play distinct roles in cortical circuit function \citep{tremblay2016gabaergic}. SST interneurons preferentially target distal dendrites of pyramidal neurons, while PVALB interneurons target perisomatic regions, leading to different effects on network dynamics. Both interneuron types express NMDA receptors with subunit compositions that confer relatively high sensitivity to ketamine's blocking action. We propose that variability in the relative density and distribution of these interneuron subtypes across individuals may contribute to the observed inter-individual variation in the principal connectivity gradient.

Several hypotheses can explain why multiple PCs emerge from a drug that primarily targets a single receptor type. Different PCs may reflect effects on different synapse types (e.g., excitatory-to-excitatory vs.\ excitatory-to-inhibitory synapses), different receptor affinities (high-affinity GluN2C/GluN2D vs.\ lower-affinity GluN2A/GluN2B), or the engagement of ketamine's multiple secondary molecular targets beyond NMDA receptors. Future studies combining pharmacological specificity with multi-dimensional neural analysis will be needed to adjudicate among these possibilities.

The demonstration that behavioral PCs map onto distinct neural connectivity gradients at the single-subject level has practical implications for clinical applications. If individual patients' neural responses to ketamine can be decomposed into interpretable dimensions that predict symptom profiles, this information could inform personalized treatment strategies, including dose optimization, combination therapies, or identification of patients likely to respond.

\subsection{Limitations}

Several limitations should be noted. First, the sample size of 40 participants, while adequate for PCA, limits the number of PCs that can be reliably estimated and the generalizability of the findings. Second, the gene expression analysis relies on postmortem AHBA data from six donors, which may not perfectly reflect in vivo expression patterns. Third, the behavioral assessments, while comprehensive, do not capture the full range of ketamine's acute subjective effects. Fourth, the cross-sectional design does not address how the multi-dimensional neural signatures relate to subsequent antidepressant response.

\section{Conclusion}

We demonstrate that the acute neural and behavioral effects of ketamine are irreducibly multi-dimensional, with the principal neural gradient linked to SST and PVALB interneuron gene expression patterns. The multi-dimensional framework resolves contradictory findings in the literature, reveals cell-type-specific mechanisms, and supports single-subject neuro-behavioral mapping. These findings advance our understanding of how a single molecular target can produce complex, individually variable neural and behavioral effects, and they lay the groundwork for precision medicine approaches to ketamine therapy.

\section*{Acknowledgments}
We thank all participants for their time and contribution to this research.

\section*{Data Availability}
Data and analysis code are available upon reasonable request to the corresponding author.

\bibliographystyle{plainnat}
\bibliography{references}

\end{document}
